% !TEX root = ../main.tex
% (this is to make TexShop know where your parent file is)

% Template questions
%
% How have you fulfilled the brief from a technical point of view?
%
% Did you have any technical difficulties?  If so, what were they and how did you overcome them? What worked well?
%
% Did you learn any new technical skills for this project?

% Grading matrix:
%
% Assessment criteria:
% TECHNICAL: Create a convincing and effective mock-up with samples. Deliver music that complies technically to the specification required and conforms to industry best practice.
%
% Intended learning outcomes:
% Deliver music that is well produced and complies technically to the brief and conforms to industry best practice.
%
% Goal (distinction):
% Music delivered complies perfectly with all the requirements of the project. Musical production and performance is of an exceptionally high quality.

% \phantomsection
% \subsubsection*{Low end theory}
% \addcontentsline{toc}{subsection}{Low end theory}

I ran into a few interesting technical issues to address:

In Guy's first thirty seconds, there is a low drone that runs continuously [00:12--00:30]. Judging from spectrum analysis, it also includes a lower octave that I can't hear on my studio monitors or on my headphones. This takes up a lot of headroom and muddies up the mix. In continuing the track, I opted for not re-introducing this after cutting it [00:40]. There is still a sub underneath the basses, but it doesn't include the super low octave.

Another issue was that the kick drum didn't seem to cut through the mix enough, [01:20, 01:39]. The issue was that the bass and kick were both occupying the same frequency range. This is a well-known problem, so I decided to duck the bass whenever the kick (or snare) sounded. This also creates a little extra bit of movement; a net positive.

I wanted to give the kick some extra impact for the final theme statement, so I layered in a heavy EDM kick and took out the clicky top end, keeping only the lows (up to ~200Hz) [01:39]. For the top layer, I used a mix of taikos and Stormdrum 2's ``Middle Earth ensemble''. Additionally, Snapback provides a reverse layer that leads into the hit for some subtle, extra impact.
